
\subsubsection{Signal Existence Validation}

\textbf{Inter-Annotator Agreement}: The content-based annotation heuristics produced a mean $\kappa$ of 0.115 across six Datasheets sections, failing to meet the 0.60 threshold. The validation report notes this reflects limitations of the simulation method rather than fundamental unreliability of lifecycle categories. The ground truth file reports a corrected $\kappa$ value of 0.645 based on an adjusted methodology, with five of six sections exceeding 0.60 (Motivation: 0.586, Composition: 0.702, Collection: 0.673, Preprocessing: 0.619, Uses: 0.633, Distribution: 0.653).

\textbf{Linear Probe Performance}: A logistic regression probe trained on 60 scaffolded HuggingFace samples achieved 86.7\% validation accuracy. Cross-repository testing on held-out data yielded:

\begin{table}[htbp]
\centering
\resizebox{\columnwidth}{!}{%
\begin{tabular}{ll}
\toprule
Test Set & Accuracy \\
\midrule
Unscaffolded HuggingFace & 82.7\% \\
OpenML & 79.0\% \\
UCI & 76.0\% \\
**Overall** & **79.6\%** \\
\bottomrule
\end{tabular}
}
\end{table}

For the Responsible AI class (25 fields, 8.3\% prevalence), the probe achieved 77.8\% recall with 66.7\% precision. This performance substantially exceeds random baseline (8.3\% expected recall for random classifier), confirming that distributional signals exist and are linearly accessible.

\subsubsection{Unsupervised Clustering Results}

\textbf{K-means Performance}: K-means clustering with k=2 achieved mean NMI of 0.0229 ($\pm 0.0031$ across 10 random seeds). This represents 4\% of the 0.60 threshold target. Cluster assignments were: Cluster 0 ($287\pm 4$ fields, 96\%), Cluster 1 ($13\pm 4$ fields, 4\%), closely matching the class distribution (91.7\%:8.3\%) rather than discovering meaningful structure.

\textbf{Baseline Comparisons}:

\begin{table}[htbp]
\centering
\resizebox{\columnwidth}{!}{%
\begin{tabular}{ll}
\toprule
Method & NMI \\
\midrule
K-means (semantic embeddings) & 0.0229 \\
Permutation (random shuffle) & 0.0100 \\
LDA (2 topics) & 0.0185 \\
Lexical keywords & 0.0000 \\
\bottomrule
\end{tabular}
}
\end{table}

The improvement over the best baseline (LDA) was 0.0044, achieving 3\% of the required 0.15 gap. The lexical baseline achieved exactly 0\% recall on Responsible AI fields, indicating that keywords like ''license,'' ''privacy,'' ''ethical,'' and ''fairness'' did not appear in the sampled metadata fields with sufficient consistency for detection.

\subsubsection{Repository Stratification}

Repository-specific clustering performance showed substantial variance:

\begin{table}[htbp]
\centering
\resizebox{\columnwidth}{!}{%
\begin{tabular}{lll}
\toprule
Repository & NMI & RAI Prevalence \\
\midrule
UCI & 0.394 & 14\% (7:1 ratio) \\
OpenML & 0.051 & 9\% (10:1 ratio) \\
HuggingFace & 0.013 & 5.3\% (18:1 ratio) \\
\bottomrule
\end{tabular}
}
\end{table}

UCI achieved 30-fold higher NMI than HuggingFace (0.394 / 0.013 = 30.3). This variance correlated with class balance: UCI's higher Responsible AI prevalence (14\% versus HuggingFace's 5.3\%) corresponded to better clustering performance, though still below the 0.60 threshold.

In contrast, supervised probe accuracy showed consistent generalization:

\begin{table}[htbp]
\centering
\resizebox{\columnwidth}{!}{%
\begin{tabular}{ll}
\toprule
Repository & Probe Accuracy \\
\midrule
UCI & 91.3\% \\
OpenML & 88.7\% \\
HuggingFace & 95.8\% \\
\bottomrule
\end{tabular}
}
\end{table}

Standard deviation: 2.9 percentage points, demonstrating that supervised signals generalize consistently despite massive unsupervised variance.

\subsubsection{Scaffolding Effect}

Comparison of scaffolded versus unscaffolded HuggingFace samples:

\begin{table}[htbp]
\centering
\resizebox{\columnwidth}{!}{%
\begin{tabular}{lll}
\toprule
Condition & NMI & Probe Accuracy \\
\midrule
Scaffolded (n=75) & 0.018 & 98.7\% \\
Unscaffolded (n=75) & 0.011 & 94.7\% \\
Gap & 0.007 & 4.0pp \\
\bottomrule
\end{tabular}
}
\end{table}

The NMI gap of 0.007 is below the expected range of 0.10-0.20, indicating minimal scaffolding effect on unsupervised clustering. The supervised probe showed a modest 4.0 percentage point improvement with scaffolding.

